% ============================================================
% ABSTRACT
% ============================================================
\begin{abstract}
\noindent
Global Positioning System (GPS) spoofing poses a critical security threat to unmanned aerial vehicles (UAVs), as adversaries can transmit counterfeit signals to manipulate a drone's perceived position without triggering conventional failure modes. Existing countermeasures frequently depend on specialised hardware---multi-antenna arrays, military-grade encrypted signals, or high-grade inertial navigation systems---that are incompatible with the size, weight, and power constraints of commercial UAV platforms. This research presents a software-only GPS spoofing detection system for PX4-based drones that leverages cross-validation between GPS-reported position and inertial measurement unit (IMU) data to distinguish spoofing from environmental disturbances such as wind. The system employs a multi-layered detection pipeline combining eight physics-based heuristic rules with machine learning classifiers---Random Forest and one-dimensional Convolutional Neural Network---operating on 30-sample sliding windows of telemetry data acquired at 10\,Hz via the MAVLink protocol. Evaluated on a dataset of 4,207 telemetry records from 8 simulated flight runs in a PX4 Software-In-The-Loop environment with Gazebo, the Random Forest achieved 100\% accuracy on the held-out test set, while the 1D CNN exhibited perfect recall but a 100\% false positive rate on the limited normal-flight sample---a finding termed the ``CNN Paradox''. The key contribution is a physics-based GPS--IMU consistency discriminator that reliably identifies spoofing without requiring prior training on attack signatures. A graduated GPS Trust Index is proposed for real-time mitigation. The results demonstrate that real-time software-only spoofing detection is feasible, though the 3-second detection window exceeds theoretical safe-reaction times for high-speed scenarios, motivating future work on reduced-latency architectures and physical hardware validation.
\end{abstract}
